% RESEARCH_DRAFT_SHA256: 58440cbb76f3f92acd5cb9dad4ee543415148818d0ee482fa59c2505c32b2c6a
% Layout and source-reference wrappers expanded for the web renderer.
\documentclass[11pt]{article}
\usepackage[margin=25mm]{geometry}
\usepackage{fontspec,unicode-math}
\setmainfont{Latin Modern Roman}
\setmathfont{Latin Modern Math}
\usepackage{amsmath,amssymb,mathtools}
\usepackage{xurl,hyperref}
\hypersetup{colorlinks=true,urlcolor=blue,linkcolor=blue}
\setcounter{secnumdepth}{0}
\setlength{\parindent}{0pt}
\setlength{\parskip}{5pt}
\setlength{\emergencystretch}{3em}
\newcommand{\sref}[2]{\hyperlink{source:#1:#2}{[S#2]}}
\newcommand{\cataloguescope}[1]{\paragraph{Source question.}#1}
\newcommand{\floor}[1]{\left\lfloor#1\right\rfloor}
\begin{document}
\section*{Research-attempt provenance}
This standalone research attempt retains the complete original definition body below,
including its historical source and preparation notes. Those notes describe the
earlier catalogue preparation; the new attempt and its verification are identified
separately after that body. No original ID, catalogue status, or source statement is changed.

% UnsolvedMath ID 30005029; source catalogue code OWR-9790359-010
\section{Determinacy Methods in LOCAL and Borel Coloring}\label{30005029}
{\small UnsolvedMath ID 30005029 · Computer Science}
\subsection{Definitions and mathematical statement}

\paragraph{Shared notation.}Standard number systems are denoted by $\mathbb N,\mathbb Z,\mathbb Q,\mathbb R,\mathbb C$; index ranges, computational models and specialized notation follow the cited problem source and the definitions below.


Write $[\mathbb N]^{\mathbb N}$ for the infinite subsets of $\mathbb N$ with their standard Borel structure. The shift map removes the least element; the undirected shift graph $G_S$ joins $x$ to $x\setminus\{\min x\}$. A proper coloring gives distinct colors to adjacent vertices; its Borel version must be Borel measurable. The chromatic number is the least number of colors, while the Borel chromatic number restricts to Borel colorings. The Axiom of Determinacy (AD) says every length-$\omega$ perfect-information game on natural numbers has a winning strategy for one player. The LOCAL model counts synchronous graph-communication rounds.
\paragraph{Statement recorded in the supplied source.}
Can shift-graph and Galvin–Prikry methods be used in the LOCAL model? Can determinacy arguments construct functions of large Borel chromatic number, and does the Axiom of Determinacy imply that the shift graph $G_S$ has infinite chromatic number? \href{https://doi.org/10.4171/owr/2022/8}{[S2]}
\paragraph{Scope qualification.}
The determinacy question is over ZF as in the source. Ordinary and Borel chromatic numbers are distinct; the final question concerns the ordinary chromatic number under AD.
\paragraph{Source status.}
The supplied catalogue labels this record open. Its source-status metadata was checked against the supplied snapshot on 2026-10-11. This import does not independently establish current open status or mathematical correctness. The catalogue literature assessment is dated 2026-08-22.
\subsection{Short English statement}
Can infinite-game and Ramsey methods inform local algorithms, and does determinacy force the infinite-subset shift graph to need infinitely many colors?
\subsection{Sources}
\begin{itemize}
\item[S1] Ulam AI and the UnsolvedMath Contributors, UnsolvedMath: A Curated Collection of Open Mathematics Problems, record 30005029 (OWR-9790359-010). Catalogue curation and metadata: CC BY 4.0. The supplied source question is retained; definitions and the English summary are editorial additions.
\url{https://huggingface.co/datasets/ulamai/UnsolvedMath}.
\item[S2] Mini-Workshop: Descriptive Combinatorics, LOCAL Algorithms and Random Processes (2023), Oberwolfach Reports, DOI 10.4171/owr/2022/8, pp. 429--455
\url{https://doi.org/10.4171/owr/2022/8}.
\item[S3] A complexity problem for Borel graphs
\url{https://doi.org/10.1007/s00222-021-01047-z}.
\textit{Underlying Oberwolfach Reports material retains its source-specific licence, including CC BY-SA 4.0 where applicable; see THIRD\_PARTY\_NOTICES.md.}
\end{itemize}
\paragraph{Preparation and review.}
Original question attributed to the cited source and UnsolvedMath contributors. Codex (GPT-6; exact serving revision not exposed) assisted with definitions, English summaries, source inspection and formatting. No new human mathematical authorship or independent proof verification is asserted. Canonical ID 30005029 and source code OWR-9790359-010 checked against the complete supplied registry and existing repository statements; known flagged repetitions and obvious identical conjectures excluded. Full mathematical-equivalence/literature review has not been performed. Statement preparation and final review followed contribution prompts C and D; automated checks verify structure and identity, not open status.
\label{end:30005029}

\section{Literature and context check}

I read the shift-graph discussion on pages 451--453 of the primary report.
It invokes Galvin--Prikry for Borel colorings and expressly says that the
implication from AD to the Ramsey property is not established there.
The calculations below therefore do not assume that implication. They examine
an exact finite shift-coloring problem and the conditional Ramsey argument,
then identify why neither gives the original determinacy result.

\section{Research attempt: Finite shift-color bounds and the choice-sensitive infinite obstruction}

\paragraph{An exact finite combinatorial calculation.}
For $N\geq2$, let $H_N$ have vertices $(i,j)$ with $0\leq i<j<N$ and
edges between $(i,j)$ and $(j,k)$ whenever $i<j<k$.
If a proper coloring uses $q$ colors, put
$S_j=\{c(i,j):i<j\}$, including $S_0=\varnothing$.
For every $j<k$, the color $c(j,k)$ belongs to $S_k$ but not $S_j$:
any occurrence $c(i,j)=c(j,k)$ would violate the edge constraint.
Thus the $N$ subsets $S_0,\ldots,S_{N-1}$ are distinct, proving
$N\leq2^q$ and $q\geq\lceil\log_2N\rceil$.
For the matching upper bound, write the integers in
$\lceil\log_2N\rceil$ bits and color $(i,j)$ by the most significant
position at which they differ. Since $i<j$, the bit there is zero in $i$
and one in $j$. If $(i,j)$ and $(j,k)$ had the same position, the bit of
$j$ would be both one and zero, impossible. Hence
\[
 \chi(H_N)=\lceil\log_2N\rceil.
\]
This is a finite constraint on color rules for consecutive pairs of identifiers.
It is not a bounded-degree tree LOCAL lower bound: these finite graphs have
growing degree and are prefix-state constructions, not asserted to be finite
subgraphs of the infinite functional shift forest.

\paragraph{The exact Ramsey hypothesis that would suffice.}
Suppose a finite coloring of $[\mathbb N]^{\mathbb N}$ has an infinite set
$A$ on which \emph{all infinite subsets} have the same color.
Choose any infinite $x\subseteq A$. Then both $x$ and
$Sx=x\setminus\{\min x\}$ are infinite subsets of $A$, so the coloring
is not proper for $G_S$. Thus a Ramsey property for every finite coloring
would imply $\chi(G_S)=\infty$.
The primary report applies the external Galvin--Prikry theorem to Borel
colorings for this reason. That theorem is context, not re-proved here.
The same reasoning cannot be applied to arbitrary color classes under AD
without a theorem establishing the corresponding Ramsey conclusion.

\paragraph{Why an ordinary two-coloring under choice is consistent.}
Orient every edge $x\to Sx$. Every vertex has exactly one outgoing edge and
there are no directed cycles, since the minimum strictly increases at each
forward step. There is no finite undirected cycle either: on such a cycle,
the number of edges equals the number of vertices, while each vertex can
supply at most one outgoing cycle edge. Equality forces exactly one outgoing
edge at every cycle vertex, giving a directed cycle, a contradiction.
Thus each component is a tree. Under the axiom of choice select one vertex
in each component and color vertices by parity of their finite graph distance
from that chosen vertex. Unique paths in a tree make this a proper two-coloring.
This selection need not be Borel and is not an operation justified by ZF+AD.
Consequently an argument that silently chooses component roots would discard
precisely the set-theoretic issue in the source question.

\section{Verification and adversarial audit}

The finite proof includes $N=2$, where the single vertex needs one color.
In the infinite graph, the lack of finite cycles is proved from the orientation,
and the use of choice in selecting component representatives is explicit.
The Ramsey implication distinguishes one homogeneous infinite family from an
ordinary infinite subset of vertices. AD is never used as an unproved substitute
for that property.

\paragraph{Audit corrections.} Did not assume AD implies the Ramsey property. Distinguished the finite
pair graph from the infinite shift forest and marked exactly where the ordinary
two-coloring invokes component-wise choice.
\paragraph{External computation.} None; symbolic arguments and boundary checks only.
\section{Scope and remaining gaps}

The finite shift bound is known elementary material re-derived as a possible
local-rule test. The original questions concerning transfer to LOCAL,
determinacy constructions and ordinary chromatic number under ZF+AD remain
unresolved; the missing regularity/choice bridge is specified rather than assumed.

\paragraph{Integration review.} Codex (GPT-6; exact serving revision not exposed) reviewed the written arguments and expanded layout and source-reference wrappers. Repository metadata and the production web renderer were checked using Python 3.13 and JavaScript with MathJax 3.2.2. This model-assisted review is not independent human proof verification; no PDF compilation or formal proof checking is asserted.

\section{Completion Estimate}
\noindent\textbf{COMPLETION: 0\%}
The calculations below the preserved source statement investigate a concrete route,
but establish no substantiated new progress toward the entire original target.
Known special cases and failed sufficient conditions are not credited as a solution.
\section{References}
The original source references and catalogue attribution are preserved in the
Sources section above. The following URLs were actually consulted for this attempt;
retention of an additional original reference does not claim fresh consultation of it.
\begin{itemize}
\item \url{https://ems.press/content/serial-article-files/46944} (consulted 2026-10-11). 
\end{itemize}
\section{Model and contributor disclosure}
Codex GPT-6 completed written-proof review of this attempt at its declared scope; this is model-assisted review, not independent human proof checking.
OpenAI gpt-6-astra, requested reasoning configuration ultra, generated this research
attempt on 2026-10-11. The exact serving revision is not exposed. Codex GPT-6
performed orchestration and repository-format review; its exact serving revision
is also not exposed. The mathematical route, boundary checks and final adversarial
self-audit followed contribution prompts A then B in
\texttt{docs/contribution\_prompts.md}, after reading \texttt{CONTRIBUTING.md}.
Python 3.14.2 assembled the standalone documents; any mathematical computation is
specified in Verification. No theorem prover was used. No human mathematical
authorship, substantial human proof correction, or independent human proof checking
is asserted. Model-assisted review, source retrieval, and structural checks do not
establish correctness or novelty. Full private reasoning transcripts are unavailable
as reproducibility artifacts; the complete arguments, inputs and outcomes actually
used are stated here or in the companion batch audit.
\section{Final status}
Unresolved as a new research contribution. This is a concrete route investigation,
not a claim of new complete resolution or a current-status certification. A known
answer reproduced for context does not become a new contribution. The catalogue's
source-status label is preserved; this attempt status does not assert that a
source-labelled solved question is open.
\end{document}
